\section{Nyman--Beurling difference budgets and their arithmetic limits (NS-59)}
\label{sec:v158-nb-difference}

\paragraph{Status.}
This section gives exact identities, proved implications, and separately
labelled finite certificates. It supplies an unconditional asymptotic upper
bound for a complete block Gram matrix. It does not supply the lower
correlation estimate needed for Nyman--Beurling convergence. An almost
scale-invariant bound for the \emph{unprojected} difference Gram is equivalent
to Lindel\"of; every fixed logarithmic version of that bound is false.
Neither assertion is transferred to the projected Gram or the particular
residual direction without an additional argument. No priority claim is made
for these consequences of the cited classical zeta estimates.

Use the Hilbert space, $\rho_n$, $P_N$, $r_N$, $E_N$, $h$ and $S$ from
Section~\ref{sec:ns53-nb-blocks}. In this section the new block is
$N<n\leq 2N$, and its indices always retain their integer values.
For $n\geq2$ put
\begin{equation}
 \delta_n=\rho_n-\frac{n-1}{n}\rho_{n-1},\qquad
 H_k=\sum_{j=1}^k\frac1j,\qquad
 L_N=\sum_{n=N+1}^{2N}\frac{H_{n-1}}{n^2}.
 \label{eq:v158-difference-def}
\end{equation}

\begin{proposition}[An elementary adjacent-dilation budget]
\label{prop:v158-adjacent-budget}
For every integer $n\geq2$,
\begin{equation}
 \frac{H_{n-1}}{n^2}-\frac{n-1}{n^3}
 \leq \norm{\delta_n}_H^2
 \leq\frac{H_{n-1}}{n^2}.
 \label{eq:v158-adjacent-norm}
\end{equation}
In particular $\norm{\delta_n}_H^2=(\log n+O(1))/n^2$ and
$L_N\leq H_{2N-1}/(2N)$.
\end{proposition}
\begin{proof}
Write $a=n-1$ and $t=1/x$. The function
\[
 \delta_n(1/t)=\frac{a}{n}\left\lfloor\frac{t}{a}\right\rfloor
              -\left\lfloor\frac{t}{n}\right\rfloor
\]
has period $an$. On $[ka,kn)$, $1\leq k\leq a$, its value is
$1-k/n$; on $[kn,(k+1)a)$, $1\leq k\leq a-1$, its value is $-k/n$.
It vanishes on $[0,a)$. Thus, including the zero-length last interval,
\begin{align*}
 \int_0^{an}\frac{|\delta_n(1/t)|^2}{t^2}\,dt
 &=\sum_{k=1}^a\left\{
 \frac{(n-k)^2}{k a n^3}
 +\frac{k(a-k)}{(k+1)a n^3}\right\}\\
 &=\frac{H_a}{n^2}-\frac{a}{n^3}.
\end{align*}
Periodicity and the displayed values give $|\delta_n(1/t)|\leq a/n$ everywhere.
The remaining nonnegative integral is at most
$(a/n)^2\int_{an}^{\infty}t^{-2}\,dt=a/n^3$.
Finally $\sum_{N<n\leq2N}n^{-2}\leq\int_N^{2N}u^{-2}\,du=1/(2N)$.
\end{proof}

Let $T_N$ be the upper bidiagonal matrix whose column indexed by $n$
has diagonal entry $1$ and, when $n>N+1$, preceding entry $-(n-1)/n$.
The first column represents $(I-P_N)\delta_{N+1}=(I-P_N)\rho_{N+1}$;
the old $\rho_N$ term is not discarded before projection. Define
\begin{equation}
 \mathcal D_N=(\langle \delta_m,\delta_n\rangle)_{N<m,n\leq2N},\quad
 K_N=T_N^*S T_N,\quad g=T_N^*h.
 \label{eq:v158-difference-gram}
\end{equation}
Here $\mathcal D_N$ is the raw difference Gram and $K_N$ is its
orthogonally projected Gram. In coordinates,
$g_n=h_n-(n-1)h_{n-1}/n$, with $h_N=0$.

\begin{proposition}[Complete difference-block gains]
\label{prop:v158-difference-gain}
The following identities and inequalities retain all old/new cross terms:
\begin{align}
 E_N-E_{2N}&=g^*K_N^{-1}g, & 0<K_N&\leq\mathcal D_N,
 \label{eq:v158-exact-difference-gain}\\
 E_N-E_{2N}&\geq\frac{\norm{g}_2^2}{L_N},&
 E_N-E_{2N}&\geq
 \frac{\norm{g}_2^4}{g^*K_Ng}\quad(g\ne0).
 \label{eq:v158-difference-bounds}
\end{align}
Replacing the old denominator by $L_N$ while keeping the old numerator
$\norm{h}_2^2$ is not justified by these formulas.
\end{proposition}
\begin{proof}
The invertible matrix $T_N$ changes the projected spanning family, so the
first identity follows from Proposition~\ref{prop:ns53-nb-block-gain}.
Orthogonal projection is contractive, hence $K_N\leq\mathcal D_N$.
Proposition~\ref{prop:v158-adjacent-budget} bounds
$\operatorname{tr}\mathcal D_N$ by $L_N$. The first lower bound follows
from $K_N\leq L_N I$, and the second is the variational gain formula
applied to the coefficient direction $g$. If $g=0$, the gain is zero.
\end{proof}

\begin{proposition}[Mellin multiplier and an unconditional block norm]
\label{prop:v158-mellin-budget}
For complex coefficients $a=(a_n)_{N<n\leq2N}$ set
\[
 F_a(u)=\sum_{n=N+1}^{2N}\frac{a_n}{n}\mathbf1_{(n-1,n]}(u),
 \qquad q_a(v)=e^{v/2}F_a(e^v).
\]
With $\widehat q(t)=\int_{\mathbb R}q(v)e^{-itv}\,dv$,
\begin{equation}
 a^*\mathcal D_Na
 =\frac1{2\pi}\int_{\mathbb R}
       |\zeta(\tfrac12+it)|^2|\widehat q_a(t)|^2\,dt,
 \qquad
 \norm{q_a}_2^2=\sum_{N<n\leq2N}\frac{|a_n|^2}{n^2}.
 \label{eq:v158-mellin-identity}
\end{equation}
If $0<\theta<1/2$ and
$|\zeta(1/2+it)|\leq A_\theta(1+|t|)^\theta$ for all real $t$, then
\begin{equation}
 \norm{\mathcal D_N}_{\rm op}
 \leq C_\theta A_\theta^2 N^{-2+2\theta},
 \qquad \norm{K_N}_{\rm op}\leq\norm{\mathcal D_N}_{\rm op}.
 \label{eq:v158-zeta-budget}
\end{equation}
Consequently, for every $\varepsilon>0$,
\begin{equation}
 \norm{\mathcal D_N}_{\rm op}\ll_\varepsilon N^{-5/3+\varepsilon}.
 \label{eq:v158-weyl-budget}
\end{equation}
This is an analytic bound with an unspecified constant, not a numerical
certificate for a finite cutoff or a claim of the best available exponent.
\end{proposition}
\begin{proof}
For $0<\Re s<1$ the usual fractional-part Mellin identity gives
$\mathcal M\rho_n(s)=-n^{-s}\zeta(s)/s$. Also
\[
 n^{-s}-\frac{n-1}{n}(n-1)^{-s}
 =\frac{1-s}{n}\int_{n-1}^n u^{-s}\,du.
\]
On $s=1/2+it$ the factor $(1-s)/s$ has modulus one. Mellin Plancherel
therefore gives~\eqref{eq:v158-mellin-identity}, including the entire
frequency integral. The norm identity follows by changing variables
$u=e^v$.

Here is the needed inverse estimate, including the discontinuities of
$q_a$. Its support is $[\log N,\log(2N)]$, partitioned into intervals
$I_n=(\log(n-1),\log n]$ of lengths between $1/(2N)$ and $1/N$.
On each interval $q_a(v)$ is a constant times $e^{v/2}$. For
$0<\theta<1/2$,
\begin{equation}
 \int_{\mathbb R}(1+|t|)^{2\theta}|\widehat q_a(t)|^2\,dt
 \leq C_\theta N^{2\theta}\norm{q_a}_2^2.
 \label{eq:v158-inverse-sobolev}
\end{equation}
To check this, use the Fourier identity between the homogeneous
$H^\theta$ seminorm and
$\iint |q(v)-q(w)|^2|v-w|^{-1-2\theta}\,dv\,dw$.
On a single cell, $q'=q/2$ bounds the integral by a constant times its
$L^2$ mass. For different cells, including the zero exterior, use
$|q(v)-q(w)|^2\leq2(|q(v)|^2+|q(w)|^2)$.
For $v\in I=(b,c)$, integration outside $I$ gives
$((v-b)^{-2\theta}+(c-v)^{-2\theta})/(2\theta)$.
Since $|q|^2$ varies by at most $e^{|I|}\leq2$ on a cell, its integral
against these endpoint powers is at most
$C_\theta |I|^{-2\theta}\int_I|q|^2$.
The powers are integrable exactly for $2\theta<1$. Summing gives
\eqref{eq:v158-inverse-sobolev}; thus no $H^1$ regularity is assumed.
Combining with~\eqref{eq:v158-mellin-identity} proves the operator bound.
The classical Weyl estimate
$\zeta(1/2+it)\ll_\delta(1+|t|)^{1/6+\delta}$
from~\cite[Section 2]{HuxleyIvic2007} gives
\eqref{eq:v158-weyl-budget} by taking $\theta=1/6+\varepsilon/2$
for small $\varepsilon$; larger $\varepsilon$ follow at once.
\end{proof}

\begin{proposition}[A sharp raw norm would already encode Lindel\"of]
\label{prop:v158-lindelof-budget}
The assertion
\begin{equation}
 \text{for every }\varepsilon>0,\qquad
 \norm{\mathcal D_N}_{\rm op}\ll_\varepsilon N^{-2+\varepsilon}
 \quad(N=2^j)
 \label{eq:v158-near-scale-budget}
\end{equation}
is equivalent to the Lindel\"of hypothesis. Moreover, for every fixed
$B\geq0$,
\begin{equation}
 \norm{\mathcal D_N}_{\rm op}
 \ne O\left(\frac{(\log N)^B}{N^2}\right)
 \quad\text{even when }N\text{ is restricted to powers of two}.
 \label{eq:v158-no-log-budget}
\end{equation}
Both statements concern the raw Gram $\mathcal D_N$ uniformly over
\emph{all} coefficient vectors. They assert neither such an equivalence
nor such an obstruction for $K_N$ or for $g^*K_Ng$.
\end{proposition}
\begin{proof}
Lindel\"of and Proposition~\ref{prop:v158-mellin-budget} imply
\eqref{eq:v158-near-scale-budget}. For the converse we first show that,
with an absolute constant and all sufficiently large $N$,
\begin{equation}
 \int_{\tau-1}^{\tau+1}|\zeta(\tfrac12+it)|^2\,dt
 \leq C N^2\norm{\mathcal D_N}_{\rm op}
 \qquad(|\tau|\leq\sqrt N).
 \label{eq:v158-local-zeta-test}
\end{equation}
Fix a nonzero nonnegative smooth function $w$ compactly supported in
$(0,\log2)$. Its Fourier transform satisfies
$|\widehat w(t)|\geq\cos(\log2)\int w>0$ for $|t|\leq1$.
Put $v_n=\log(n-1/2)$ and choose
\[
 a_n=n e^{-v_n/2}w(v_n-\log N)
                      e^{i\tau(v_n-\log N)}.
\]
The associated $q_a$ approximates
$q(v)=w(v-\log N)e^{i\tau(v-\log N)}$ with
\[
 \norm{q_a-q}_{L^1}\leq C_w\frac{1+|\tau|}{N},
 \qquad \norm{a}_2^2\leq C_w N^2.
\]
These estimates follow cell by cell from the mean value theorem and
$|v-v_n|\leq C/N$; the constants are independent of $N$ and $\tau$.
For $|\tau|\leq\sqrt N$, the Fourier error tends uniformly to zero.
Hence $|\widehat q_a(t)|$ is bounded below by a fixed positive constant
on $[\tau-1,\tau+1]$ for all sufficiently large $N$.
Equation~\eqref{eq:v158-mellin-identity} proves
\eqref{eq:v158-local-zeta-test}. Complex coefficients are harmless:
$\mathcal D_N$ is real symmetric, and its operator norm is the same
over real or complex vectors.

The classical short-interval inequality
\cite[Eq. (2.1)]{HuxleyIvic2007} is
\begin{equation}
 |\zeta(\tfrac12+iT)|^2\ll\log T
 \left(1+\int_{T-\log^2T}^{T+\log^2T}
                      |\zeta(\tfrac12+it)|^2\,dt\right)
 \quad(T\geq T_0).
 \label{eq:v158-short-interval}
\end{equation}
Choose a power of two $N$ with $4T^2\leq N<8T^2$.
For large $T$ the integration interval can be covered by
$O(\log^2T)$ intervals in~\eqref{eq:v158-local-zeta-test}.
Consequently
\begin{equation}
 |\zeta(\tfrac12+iT)|^2
 \ll\log T+(\log T)^3 N^2\norm{\mathcal D_N}_{\rm op}.
 \label{eq:v158-budget-to-zeta}
\end{equation}
Assertion~\eqref{eq:v158-near-scale-budget}, with arbitrarily small
exponents, now implies Lindel\"of.
If the big-$O$ estimate excluded in~\eqref{eq:v158-no-log-budget}
held, the same formula would give
$|\zeta(1/2+iT)|^2\ll(\log T)^{B+3}$.
This contradicts the unconditional large values in
\cite[Theorem 1]{Soundararajan2008}.
\end{proof}

\paragraph{Certified finite comparison; no extrapolation.}
The exact rational-cotangent formula for the Gram entries is taken from
\cite[Proposition 89]{BaezDuarteAutocorrelation2003}. With
$V(p,q)=\sum_{k=1}^{q-1}\{kp/q\}\cot(\pi k/q)$, $d=(m,n)$ and
$\kappa=\log(2\pi)-\gamma$, its scaling to the current conventions is
\begin{equation}
 G_{mn}=\frac\kappa2\left(\frac1m+\frac1n\right)
 +\frac{n-m}{2mn}\log\frac mn
 -\frac{\pi d}{2mn}\{V(m/d,n/d)+V(n/d,m/d)\}.
 \label{eq:v158-cot-gram}
\end{equation}
All entries, cross terms and solves in \texttt{evidence/v158/nb\_certify.py}
are evaluated as Arb balls, independently at 256 and 384 bits, for the
same matrices through index 512. The program also checks
\eqref{eq:v158-adjacent-norm} for $2\leq n\leq512$ as a finite replay
of the analytic estimate. At the displayed blocks, the following numbers
are outward-rounded enclosing intervals; all ratios are relative to $E_N$:
\begin{center}
\begin{tabular}{rcccc}
\toprule
$N$ & actual gain & raw-trace bound & difference budget & difference direction\\
\midrule
128 & $(.14752,.14753)$ & $(.0003503,.0003505)$
    & $(.0004286,.0004289)$ & $(.08380,.08381)$\\
256 & $(.10289,.10290)$ & $(.0002135,.0002137)$
    & $(.0001328,.0001331)$ & $(.05955,.05956)$\\
\bottomrule
\end{tabular}
\end{center}
The raw-trace bound is less than $1/400$ of the actual gain at both
blocks. At $N=256$ the smaller adjacent-difference denominator produces
an even smaller trace bound after transforming the numerator correctly.
The full directional quotient performs better on these two finite blocks;
this is not an asymptotic lower bound. The floating scan through index
2048 is separately labelled a numerical illustration in
\texttt{evidence/diag\_ns59\_nb\_budget/}.

\paragraph{What remains missing.}
The unconditional bound~\eqref{eq:v158-weyl-budget} supplies a denominator,
so it gives $E_N-E_{2N}\geq C_\varepsilon^{-1}
N^{5/3-\varepsilon}\norm{g}_2^2$ for a suitable constant.
To use it for convergence one still needs, for example, a cofinal lower
bound $\norm{g}_2^2\geq C_\varepsilon N^{-5/3+\varepsilon}
\alpha_j E_N$ at $N=2^j$, with $0\leq\alpha_j<1$ and
$\sum_j\alpha_j=\infty$. No such estimate is proved here.
The nonsummability of the \emph{exact} relative gain is simply the known
convergence criterion, not an independent arithmetic input.
Uniform logarithmic control of $\mathcal D_N$ is ruled out, while control
restricted to the actual residual direction remains open. The raw RBC
input of Section~\ref{sec:ns53-nb-blocks} is neither disproved by the
finite losses nor established by the new denominator estimates.
RH, G2 and the uniform Weil floor remain open.
